\chapter{Shaders, and why your cubes look flat}

Chapter 11 told you that plain \lstinline|DrawCube| does no lighting, and left it
there. This is the chapter that explains what to do about it.

\section{What a shader actually is}

A shader is a small program that runs \textbf{on the GPU}, in parallel, an
absurd number of times per frame. You write it in GLSL, a C-like language, and
you get two of them:

\begin{center}
\begin{tikzpicture}[
  node distance=7mm,
  every node/.style={font=\sffamily\small},
  bx/.style={draw=codeframe,fill=rlight,rounded corners=2pt,minimum width=4.6cm,minimum height=9mm,align=center},
  arrow/.style={-{Stealth[length=5pt]},rgrey}
]
\node[bx] (verts) {your vertices};
\node[bx,right=14mm of verts,fill=white,draw=rblue!50] (vs) {\textbf{VERTEX} shader\\[-1pt]{\footnotesize once per corner}};
\node[bx,below=of vs,fill=white,draw=rorange!60] (fs) {\textbf{FRAGMENT} shader\\[-1pt]{\footnotesize once per pixel}};
\node[bx,left=14mm of fs] (screen) {pixels on screen};
\draw[arrow] (verts) -- (vs);
\draw[arrow] (vs) -- node[right,font=\sffamily\tiny,text=rgrey] {where things land} (fs);
\draw[arrow] (fs) -- node[above,font=\sffamily\tiny,text=rgrey] {what colour} (screen);
\end{tikzpicture}
\end{center}

\begin{itemize}[leftmargin=*,itemsep=3pt]
  \item The \textbf{vertex} shader runs once per corner of your geometry and
        decides \emph{where} it lands on screen.
  \item The \textbf{fragment} shader runs once per pixel and decides \emph{what
        colour} that pixel is.
\end{itemize}

At 1080p a fullscreen effect runs your fragment shader two million times, sixty
times a second. That is why it must be tiny and why it cannot talk back to your
C code.

\keyidea{raylib ships a default pair that just draws things normally. You almost
always write only the fragment shader and pass \lstinline|NULL| for the vertex
one. Lighting, colour grading, outlines, water, fog --- all fragment shaders.}

\section{The three steps, which never change}

\begin{lstlisting}
// 1. LOAD, once, outside the loop
Shader shader = LoadShader(NULL, "glow.fs");          // from files
Shader shader = LoadShaderFromMemory(NULL, fsCode);   // from a string

// 2. FIND THE SLOTS, once: a lookup by name into the compiled program
int locTime = GetShaderLocation(shader, "uTime");

// 3. FEED AND USE, every frame
SetShaderValue(shader, locTime, &time, SHADER_UNIFORM_FLOAT);

BeginShaderMode(shader);
    DrawTexture(tex, 0, 0, WHITE);     // drawn THROUGH the shader
EndShaderMode();
\end{lstlisting}

That \lstinline|Begin|/\lstinline|End| pair is the same contract as the cameras:
inside it, the shader applies; outside, normal drawing. Your HUD goes outside,
for the same reason as always.

\section{Uniforms: the only way to talk to the GPU}

A \emph{uniform} is a variable you send from C into the shader. It is called
uniform because it has the same value for every pixel in that draw call.

\begin{center}
\begin{tabular}{@{}ll@{}}
\toprule
\textbf{C side} & \textbf{GLSL side} \\
\midrule
\lstinline|SHADER_UNIFORM_FLOAT| & \lstinline|uniform float uTime;| \\
\lstinline|SHADER_UNIFORM_VEC2| & \lstinline|uniform vec2 uLight;| \\
\lstinline|SHADER_UNIFORM_VEC3| & \lstinline|uniform vec3 uColour;| \\
\lstinline|SHADER_UNIFORM_INT| & \lstinline|uniform int uMode;| \\
\lstinline|SetShaderValueTexture| & \lstinline|uniform sampler2D uSecond;| \\
\bottomrule
\end{tabular}
\end{center}

\gotcha{\lstinline|GetShaderLocation| returns $-1$ if the name is not found ---
and GLSL compilers \emph{delete unused uniforms}, so a uniform you declared but
never read in the shader body will not be there. A silent $-1$ that makes
\lstinline|SetShaderValue| do nothing is the classic first shader bug.}

\section{A fragment shader, line by line}

\begin{lstlisting}[style=plain]
#version 330

in vec2 fragTexCoord;             // from the vertex shader: where we are, 0..1
in vec4 fragColor;                // the tint passed to DrawTexture

uniform sampler2D texture0;       // the texture raylib binds automatically
uniform vec4 colDiffuse;          // raylib's diffuse colour

out vec4 finalColor;              // what we must write

void main()
{
    vec4 texel = texture(texture0, fragTexCoord)*colDiffuse*fragColor;
    float grey = dot(texel.rgb, vec3(0.299, 0.587, 0.114));
    finalColor = vec4(grey, grey, grey, texel.a);
}
\end{lstlisting}

Four things to notice:

\begin{enumerate}[leftmargin=*,itemsep=3pt]
  \item \lstinline|#version 330| must be the very first thing in the file or
        string. Not a comment, not a blank line. On desktop raylib uses GLSL 330.
  \item \lstinline|texture0|, \lstinline|colDiffuse|, \lstinline|fragTexCoord|
        and \lstinline|fragColor| are names raylib fills in for you. Spell them
        exactly or you get nothing.
  \item \lstinline|vec4| is four floats, and colours here run \textbf{0.0 to
        1.0}, not 0 to 255.
  \item The last line of the function must write \lstinline|finalColor|. That is
        the shader's entire output.
\end{enumerate}

Those weights \lstinline|(0.299, 0.587, 0.114)| are how bright each channel looks
to a human eye --- green far more than blue. Averaging the three instead gives a
muddy grey, which is why this particular triple turns up everywhere.

\section{Errors show up at run time, not compile time}

Your C compiles fine whatever nonsense is in the shader string: the GPU compiler
only runs when \lstinline|LoadShader| executes. Watch the terminal:

\begin{lstlisting}[style=plain]
INFO: SHADER: [ID 5] Fragment shader compiled successfully
INFO: SHADER: [ID 6] Program shader loaded successfully
\end{lstlisting}

If instead you see \lstinline|SHADER: [ID 5] Failed to compile fragment shader|,
the line number in the message refers to the \emph{shader} source. A failed
shader falls back to the default one, so your scene keeps rendering normally and
the only sign of trouble is that your effect does nothing.

\section{Lighting, honestly}

Real lighting is a fragment shader that receives the surface normal and the light
positions, and computes how much each light hits each pixel. raylib does not
build that in --- the \file{examples/shaders/} folder has it, in
\file{shaders\_basic\_lighting.c}, as about 80 lines of GLSL.

Before you go there, be aware of the cheaper options from chapter 11: wireframes
over solid shapes, colour variation per object, and a distinct floor and ceiling.
They cost nothing, they work everywhere, and plenty of good-looking games never
do anything more.

\keyidea{Learn shaders when flat colour is actually holding your game back, not
before. It is the deepest rabbit hole in graphics and the easiest place to spend
three weeks and ship nothing.}

\tryit{Module 3 has three fragment shaders written as C strings --- greyscale, a
wave distortion driven by a \lstinline|uTime| uniform, and a spotlight that
follows your mouse. The HUD is drawn outside the pair so you can see what the
shader does and does not touch.}{./cheatsheet/build.sh 3}
